% Flat single-file source generated from index.tex
\documentclass[conference,letterpaper]{IEEEtran}
\usepackage[T1]{fontenc}
\usepackage{amsmath}
\usepackage{amsfonts}
\usepackage{amssymb}
\usepackage{cite}

\usepackage[pdftex]{graphicx}

\usepackage[caption=false,font=footnotesize]{subfig}

\usepackage{array}
\usepackage{booktabs}
\newcolumntype{L}[1]{>{\raggedright\arraybackslash}p{#1}}

\usepackage[hidelinks,hypertexnames=false]{hyperref}

\graphicspath{{./}{Figures/Plots_And_Graphs/}}
\DeclareGraphicsExtensions{.pdf,.jpeg,.png}

\hyphenation{op-tical net-works semi-conduc-tor}

\begin{document}

\title{Benchmarking Dynamic Frequency Estimators for Low-Inertia IBR Grids: A Latency--Robustness Trade-off Analysis}
\author{
\IEEEauthorblockN{Jorge Luis Mayorga Taborda,\; Yessica Africano Rodriguez,\; Fernando Jimenez}
\IEEEauthorblockA{
Department of Electrical and Electronics Engineering\\
Universidad de los Andes\\
Bogot\'a, Colombia\\
jl.mayorga236@uniandes.edu.co, ya.africano1030@uniandes.edu.co,\\
fjimenez@uniandes.edu.co
}
}
\markboth{5th International Conference on Smart Grid Synchronized Measurements and Analytics (SGSMA 2026), Santiago, Chile, June 1--4, 2026}%
{Mayorga \MakeLowercase{\textit{et al.}}: Benchmarking Dynamic Frequency Estimators for Low-Inertia IBR Grids}

\hypersetup{
  pdftitle={Benchmarking Dynamic Frequency Estimators for Low-Inertia IBR Grids: A Latency-Robustness Trade-off Analysis},
  pdfauthor={Jorge Luis Mayorga Taborda, Yessica Africano Rodriguez, Fernando Jimenez},
  pdfsubject={SGSMA 2026 conference paper},
  pdfkeywords={frequency estimation, low-inertia power systems, inverter-based resources, Monte Carlo benchmark, relay protection, trip-risk assessment}
}

\maketitle

\begin{abstract}
Frequency estimation in inverter-based resource (IBR) grids must contend with
fast ramps, phase jumps, harmonics, ringdown, and out-of-band components that
can overlap within a single local event. This paper presents a comparative
benchmark of 16 dynamic frequency estimators across 34 test scenarios grouped
into 10 disturbance families in a local relay measurement setting. The test set
combines compliance-style dynamic tests, including magnitude steps, ramps,
modulation, and phase jumps, with IBR-specific stress conditions such as
harmonic injection, power-imbalance ringdown, out-of-band interference, and
composite multi-event sequences. Each estimator is tuned via scenario-specific
parameter tuning and evaluated using RMSE, peak error, trip-risk duration,
settling time, and CPU cost. Rankings vary by disturbance family and
performance metric: SOGI-PLL achieves the lowest error on ramps, EKF on
magnitude steps and out-of-band interference, and RA-EKF on modulation, while
the composite IBR event splits RMSE leadership (PLL) from trip-risk leadership
(IpDFT). The hardest cases are the multi-event IBR sequence, $60^\circ$ phase
jumps, and noisy ringdown variants. The benchmark shows that compliance-style
dynamic tests provide necessary but insufficient validation for
protection-oriented frequency measurement in IBR-dominated grids.
\end{abstract}
\begin{IEEEkeywords}
Frequency estimation, inverter-based resources, low-inertia systems, relay
protection, signal processing, trip-risk assessment.
\end{IEEEkeywords}

\section{Introduction}\label{sec:introduction}

Frequency estimation affects relaying, under-frequency load shedding,
disturbance monitoring, and converter supervision in low-inertia grids. In
inverter-based-resource (IBR) systems, the estimation problem is not limited
to slow electromechanical drift. Converter-dominated dynamics compress the
disturbance timescale and can combine fast RoCoF segments, abrupt phase jumps,
harmonics, out-of-band components, and short ringdown intervals within the
same protection decision window~\cite{Milano2018LowInertia,Pinheiro2025,Ponce2025}.
Recent IBR studies also show that converter interactions can reshape the
observable frequency response itself~\cite{Bernal2025,Zoghby2024}.

These effects stress estimators through different failure modes. Phase
discontinuities can be interpreted by loop-based methods as large frequency
excursions or cycle slips. Sustained ramps stress the tracking bandwidth of
multi-cycle spectral methods. Harmonic and interharmonic content increases
spectral leakage and can bias estimators whose signal model assumes a single
dominant sinusoid. Ringdown adds a transient-recovery problem in which average
error, peak error, and trip-risk duration need not favor the same method.

The estimator families used in practice reflect this latency--robustness
tension. Windowed spectral methods such as IpDFT and Taylor--Fourier variants
reject steady harmonic content but pay observation delay~\cite{Belega2014,PlatasGarza2011}.
PLL and SOGI-based trackers are light enough for embedded deployment but can
lose phase alignment under abrupt discontinuities~\cite{Golestan2017,Wright2019}.
Kalman-filter variants can track fast transients when the state model remains
valid~\cite{Ferrero2020EKF,Chamorro2021}. Reduced-order and data-driven methods
add further cost--accuracy trade-offs, but their behavior under
relay-oriented composite stress remains less standardized~\cite{Netto2018KoopmanKF,Ting2022}.

This paper frames the contribution as a benchmark study rather than as a
single-estimator paper. The benchmark extends standard dynamic-test families
with IBR-specific harmonic, out-of-band, ringdown, and multi-event stress in a
local, non-time-synchronized setting. The study compares 16 estimators over 34
scenarios using RMSE, peak error, trip-risk duration, settling behavior, and
relative CPU cost. The contribution is therefore the disturbance-family
comparison protocol and the resulting map of trade-offs; no new estimator
family is proposed.

\section{Related Work and Algorithmic Landscape}
\label{sec:related}

This paper addresses local frequency measurement for relay and monitoring
functions, not PMU time-stamp quality. IEC/IEEE 60255-118-1 and the earlier
IEEE C37.118.1 define dynamic tests around ramps, magnitude steps, modulation,
and selected phase events~\cite{IEEE_C37_118,IEC_IEEE_60255_118_1_2018}. These
tests are necessary reference cases, but they do not by themselves describe
how estimators behave when phase discontinuities, harmonics, ramps, and
ringdown occur in the same local record.

Windowed spectral methods, PLL variants, Kalman filters, and reduced-order or
data-driven estimators each occupy a different point on the latency--robustness
plane. IpDFT and TFT are not interchangeable. IpDFT uses a stationary spectral
model with a Hanning window and parabolic interpolation around the dominant bin
~\cite{Belega2014}. TFT uses a second-order Taylor-series expansion of the
fundamental component only, also under Hanning weighting, which gives explicit
phasor-derivative tracking at the cost of higher noise sensitivity
~\cite{PlatasGarza2011}.
PLL variants are computationally light but can misread abrupt phase
steps~\cite{Golestan2017,Wright2019}. EKF- and UKF-based approaches can improve
transient tracking when the state model matches the event, and RoCoF
augmentation in EKF-type estimators has been studied previously~\cite{Ferrero2020EKF,Chamorro2021}.
Reduced-order and data-driven methods add further cost trade-offs~\cite{Netto2018KoopmanKF,Ting2022}.

Prior low-inertia and IBR studies show why this gap matters~\cite{Milano2018LowInertia,Pinheiro2025,Ponce2025,Bernal2025,Zoghby2024}.
The remaining need is a comparison that reports standard dynamic tests,
composite IBR stress, trip-risk, and compute cost in the same frame. The
benchmark used here keeps the standard disturbance families, adds IBR-specific
harmonic, ringdown, out-of-band, and multi-event stress, and reports
family-level summaries rather than a single global winner. This framing
targets protection-oriented estimator selection, not synchrophasor compliance
certification.

\section{Methodology and Test Framework}
\label{sec:methodology}

The benchmark compares 16 frequency estimators on a common dynamic test set in
a local, non-time-synchronized measurement setting, using waveform generation
at $1$\,MHz and estimator outputs at $10$\,kHz. The scenario set extends the
same protection-oriented stress modes used in standard dynamic
testing--magnitude steps, ramps, modulation, and phase jumps--with
IBR-specific harmonic, out-of-band, ringdown, and multi-event stress. This
expansion does not change the measurement setting: all methods are evaluated as
local relay-class frequency trackers from the same decimated $10$\,kHz data
stream.

The estimator pool spans six practical groups: zero-crossing and energy
baselines (ZCD, TKEO), windowed spectral methods (IpDFT, TFT), PLL/SOGI-based
trackers (PLL, SOGI-PLL, SOGI-FLL, Type-3 SOGI-PLL), Kalman-filter methods
(LKF2, EKF, UKF, RA-EKF), parametric/subspace methods (Prony, ESPRIT), and
reduced-order or adaptive/data-driven designs (RLS, Koopman/RK-DPMU). This mix
is intended to expose trade-offs relevant to protection and monitoring:
accuracy, transient risk, settling behavior, and computational burden. Within
the Kalman-filter group, EKF and UKF use a single-tone state over phase,
frequency, and amplitude, whereas RA-EKF adds an explicit RoCoF state,
innovation-scaled process noise, and a soft phase reset under large
innovations. RoCoF augmentation in EKF-type filters has prior literature
antecedents~\cite{Ferrero2020EKF,Chamorro2021}; the benchmark treats RA-EKF as
one Kalman-family variant, not as a new estimator family.

The study contains 34 benchmark scenarios. Rather than read the benchmark as a
flat scenario list, we group the scenarios by disturbance family.
Table~\ref{tab:benchmark_scope} summarizes the resulting ten-family structure,
and Fig.~\ref{fig:scenario_overview} shows representative ground-truth signals
for six of those families.

\begin{table}[t]
  \centering
  \caption{Benchmark scope. Counts refer to benchmark scenarios. OOB denotes
    out-of-band interference. The phase-jump family includes $20^\circ$ and
    $60^\circ$ re-lock tests.}
  \label{tab:benchmark_scope}
  \scriptsize
  \setlength{\tabcolsep}{3pt}
  \renewcommand{\arraystretch}{1.08}
  \begin{tabular}{L{1.7cm} c L{3.95cm}}
    \toprule
    Family & Count & Stress tested \\
    \midrule
    IEEE ramps & 9 & Sustained RoCoF from 0.25 to 20\,Hz/s \\
    Mag.\ steps & 7 & 1\% to 50\% magnitude changes \\
    Modulation & 3 & Nominal, AM, and FM tracking \\
    Phase jumps & 3 & $20^\circ$ and $60^\circ$ phase-jump re-lock tests \\
    OOB interference & 1 & Spectral rejection under off-nominal content \\
    Single-tone & 1 & Deterministic steady-state control case \\
    IBR harmonics & 3 & Small, medium, and large harmonic stress \\
    IBR ringdown & 5 & Power-imbalance ringdown with noise sweep \\
    IBR multi-event & 1 & Composite phase jump, ramp, harmonics, and ringdown \\
    Freq.\ step & 1 & Isolated frequency-step control case \\
    \bottomrule
  \end{tabular}
\end{table}

The family structure reflects distinct physical stress modes. Ramps probe
sustained RoCoF tracking. Magnitude steps and modulation probe
frequency-amplitude decoupling. Phase jumps probe re-lock behavior and
transient overshoot. Out-of-band interference and IBR harmonics probe spectral
leakage. Ringdown and multi-event cases combine these effects under changing
noise and transient recovery, which makes them the strongest stress tests in
the benchmark. Two records drive the sharpest separations: a NERC-type
$+60^\circ$ phase-jump case at $t=0.7$\,s with 4\% 5th, 2\% 7th, and 0.5\%
32.5-Hz interharmonic content, and an IBR multi-event case with a 0.5-pu
sag/swell, a $+45^\circ$ phase jump, a $-2$\,Hz/s RoCoF ramp clamped at
58\,Hz, decaying DC, AM/PM oscillations, a full IBR harmonic spectrum, and
compound noise. Figure~\ref{fig:scenario_overview} illustrates that
difference directly: some benchmark families are isolated compliance-style
events, whereas others stack abrupt discontinuities, ramps, and oscillatory
recovery inside the same record.

\begin{figure*}[t]
  \centering
  \includegraphics[width=0.98\textwidth,height=0.45\textheight,keepaspectratio]{Fig1_Scenarios_Final.pdf}
  \caption{Representative benchmark scenarios. Each row shows the voltage
    waveform and true frequency trajectory for six disturbance families:
    10\% magnitude step, frequency ramp, modulation, $60^\circ$ phase jump,
    IBR multi-event sequence, and IBR power-imbalance ringdown. Dashed lines
    mark disturbance onset. The figure illustrates why the benchmark is read by
    disturbance family: some cases isolate standard dynamic stress, whereas
    others combine discontinuities, ramps, harmonics, and oscillatory recovery
    in one record. Panel (E) shows the dominant frequency profile of the
    multi-event sequence; the full scenario also includes a 0.5-pu sag/swell,
    AM/PM modulation, a full IBR harmonic spectrum, and compound noise.}
  \label{fig:scenario_overview}
\end{figure*}

All estimators are evaluated as local relay-class frequency trackers from the
same decimated $10$\,kHz data stream; the study does not assess PMU
time-stamping or synchrophasor compliance. Each estimator--scenario pair is
tuned via scenario-wise parameter search across 100 trials and then evaluated
over 60 Monte Carlo runs. The reported values therefore approximate
best-in-range performance within the tested search space, not untuned default
settings. The study contains 544
estimator--scenario summaries derived from 32,640 Monte Carlo records. Let
$e[k]=\hat{f}[k]-f_{\mathrm{true}}[k]$. We report
$\mathrm{RMSE}=\sqrt{N^{-1}\sum_k e[k]^2}$, maximum peak error
$\max_k |e[k]|$, trip-risk duration
$T_{\mathrm{trip}}=\Delta t\sum_k \mathbf{1}\{|e[k]|>0.5\,\mathrm{Hz}\}$, settling time,
and CPU time per sample. The 0.5-Hz threshold is used as a
protection-oriented risk proxy rather than as a universal trip criterion. CPU
time is interpreted as a relative computational-cost indicator measured in one
software environment, not as hardware latency. For any metric $m$, the family
summary is $\bar{m}_F=|F|^{-1}\sum_{s\in F}\bar{m}_s$, where $\bar{m}_s$ is
the scenario mean over the 60 realizations. This prevents dense families from
dominating the comparison. Unless stated otherwise, frequency errors are
reported in mHz in the tables and discussion.

\section{Simulation Results and Benchmarking}
\label{sec:results}

\begin{table*}[t]
  \centering
  \caption{Core disturbance-family comparison for the benchmark study.
    RMSE values are reported in mHz and computed as family means of
    scenario means. The uncertainty term is the mean within-scenario Monte
    Carlo standard deviation for the winning estimator. $\Delta$ is the RMSE gap
    between the winner and the runner-up.}
  \label{tab:family_summary}
  \scriptsize
  \setlength{\tabcolsep}{2.8pt}
  \renewcommand{\arraystretch}{1.1}
  \begin{tabular}{l c l c l c l}
    \toprule
    Family & Count & RMSE winner & RMSE [mHz] & Runner-up & $\Delta$ [mHz] & Lowest $T_{\mathrm{trip}}$ [s] \\
    \midrule
    IEEE ramps & 9 & SOGI-PLL & $19.4 \pm 6.4$ & Koopman & 2.7 & 0 (multiple) \\
    Mag.\ steps & 7 & EKF & $0.388 \pm 0.143$ & RA-EKF & 1.45 & 0 (multiple) \\
    Modulation & 3 & RA-EKF & $17.8 \pm 21.4$ & Koopman & 13.7 & 0 (multiple) \\
    Phase jumps & 3 & SOGI-FLL & $174 \pm 67.5$ & PLL & 4.03 & ZCD: 0.0194 \\
    OOB interf. & 1 & EKF & $13.6 \pm 11.1$ & RA-EKF & 7.35 & EKF, RA-EKF: 0 \\
    IBR harmonics & 3 & EKF & $70.5 \pm 27.3$ & TFT & 4.13 & TFT, Prony: 0.00947 \\
    IBR ringdown & 5 & UKF & $313 \pm 23.2$ & EKF & 6.67 & ESPRIT: 0.0989 \\
    IBR multi-event & 1 & PLL & $173.5 \pm 66.7$ & UKF & 2.74 & IpDFT: 0.0439 \\
    \bottomrule
  \end{tabular}
\end{table*}

\subsection{Family-Level Accuracy and Trip-Risk}

Table~\ref{tab:family_summary} summarizes family-level RMSE and trip-risk
results. The ranking changes across disturbance families. SOGI-PLL gives the
lowest mean RMSE on IEEE ramp cases (19.4\,mHz). EKF leads magnitude steps
(0.388\,mHz) and out-of-band interference (13.6\,mHz), and RA-EKF leads the
modulation family (17.8\,mHz). Zero-trip-risk ties are common in ramps,
magnitude steps, and modulation, so those families are separated mainly by
accuracy rather than by protection-oriented risk.

Across the 34 benchmark scenarios, nine different estimators win at least one
RMSE scenario, and no estimator wins more than 10 of them. The
magnitude-step family is the clearest exception: EKF wins all seven
magnitude-step scenarios and the family mean as well. This is also the family
with the clearest aggregate margin: EKF leads RA-EKF by 1.45\,mHz, while the
winner's mean within-scenario spread is 0.143\,mHz.

Most other families are tighter. In ramps, modulation, phase jumps, harmonics,
ringdown, and out-of-band interference, the family-level margin is smaller than
or comparable to the same spread. For those families, the winner should be
interpreted as a near-front method rather than as a decisive leader.

Trip-risk is less discriminating in the clean standard cases. In 23 of the 34
scenarios, the minimum trip-risk is zero, often shared by 10 to 13 estimators.
Phase jumps are the main exception. SOGI-FLL gives the lowest mean RMSE across
the phase-jump family (174\,mHz), but ZCD gives the lowest mean trip-risk
($0.0194$\,s). Even inside the standard test set, RMSE and trip-risk are not
interchangeable objectives.

\subsection{Composite IBR Stress Cases}

The largest ranking changes appear under composite stress. The hardest
benchmark scenarios, measured by median RMSE across estimators, are the IBR
multi-event sequence (826\,mHz), the two $60^\circ$ phase-jump cases
(675\,mHz and 637\,mHz), and the noisy IBR ringdown variants
(474--528\,mHz). Median RMSE in these cases exceeds the magnitude-step family
by roughly three orders of magnitude.

The IBR multi-event sequence is the clearest example. PLL gives the lowest
RMSE at 173.5\,mHz. UKF trails by only 2.74\,mHz on the family mean. EKF
remains close at 176.4\,mHz while reducing trip-risk to $0.0874$\,s. IpDFT
gives the lowest trip-risk ($0.0439$\,s) but at the cost of a much larger peak
error (2550\,mHz). The ranking therefore depends on the deployment objective:
minimizing average error, minimizing time above the trip threshold, or bounding
transient severity.

The same metric split appears in the other IBR-specific families. EKF gives the
lowest mean RMSE across the harmonic family (70.5\,mHz), but TFT and Prony
give the lowest mean trip-risk there ($0.00947$\,s). In the ringdown family,
UKF gives the lowest mean RMSE (313\,mHz), whereas ESPRIT gives the lowest
mean trip-risk ($0.0989$\,s) at a much higher computational cost.

Among the 12 scenarios with a unique trip-risk winner, nine differ from the
RMSE leader. Those disagreements concentrate in the $60^\circ$ phase-jump
cases and in the harmonic, ringdown, and multi-event IBR cases. This metric
split is why a single scalar ranking is insufficient for
protection-oriented selection.

\subsection{Global Aggregation and Metric Dependence}

Family-level reading remains the primary analysis, but global summaries still
help as secondary evidence. Their interpretation depends on the aggregation
rule. Mean-rank summaries across the 34 benchmark scenarios favor TFT on RMSE
($4.38$) and peak error ($4.32$), with RA-EKF close behind on both ($4.62$
and $4.59$). Simple averaging over the 34 scenarios instead favors RA-EKF and
EKF, with mean RMSE values of 116\,mHz and 123\,mHz, followed by UKF at
131\,mHz. The difference is informative: mean rank rewards repeated near-front
finishes, whereas raw averaging is driven more strongly by the hardest
families.

Consistency statistics make that distinction concrete. EKF is the most frequent
accuracy leader, with 10 outright RMSE wins, 19 RMSE top-three finishes, and
20 peak-error top-three finishes across the 34 scenarios. TFT wins only two
RMSE scenarios, yet it reaches the RMSE top three in 16 scenarios and the
peak-error top three in 17. Trip-risk favors a different group again: IpDFT
and ESPRIT enter the trip-risk top three in 25 and 24 scenarios, respectively.
Global summaries are therefore screening tools rather than replacements for
family-level evidence.

\begin{table}[t]
  \centering
  \caption{Compact global screening indicators across the 34 benchmark
    scenarios. These statistics summarize different aggregation rules and are
    intended as screening aids rather than as a universal ranking.}
  \label{tab:global_screening}
  \scriptsize
  \setlength{\tabcolsep}{3pt}
  \renewcommand{\arraystretch}{1.08}
  \begin{tabular}{L{1.85cm} L{2.1cm} c}
    \toprule
    Global view & Best estimator & Value \\
    \midrule
    Mean RMSE & RA-EKF & 116\,mHz \\
    Mean-rank RMSE & TFT & 4.38 \\
    RMSE wins & EKF & 10 / 34 \\
    RMSE top-three & EKF & 19 / 34 \\
    Peak-error top-three & EKF & 20 / 34 \\
    Trip-risk top-three & IpDFT & 25 / 34 \\
    \bottomrule
  \end{tabular}
\end{table}

Table~\ref{tab:global_screening} makes the aggregation split explicit. The
rank-based view favors repeated near-front finishes, the mean-RMSE view is more
sensitive to the hardest families, and the hit-count view rewards broad
coverage. No estimator leads all three views at once. Across the 544
estimator--scenario summaries, RMSE correlates strongly with trip-risk
($\rho=0.890$) and settling time ($\rho=0.867$), but only weakly and
negatively with CPU cost ($\rho=-0.202$).

\subsection{Computational Cost}

CPU cost changes the practical ranking again. Figure~\ref{fig:dashboard_mega}
gathers representative time-domain cases together with the global summaries.
EKF and RA-EKF occupy the low-error middle-cost region with mean CPU costs of
$0.813\,\mu$s and $1.61\,\mu$s per sample. SOGI-PLL and SOGI-FLL remain much
cheaper ($0.0867\,\mu$s and $0.0348\,\mu$s) and therefore stay relevant for
embedded protection even when they lose some accuracy. Window-based methods
such as TFT and IpDFT improve trip-risk in some composite cases at about
$3$--$4\,\mu$s per sample. Koopman and ESPRIT occupy a much higher-cost
region, with mean costs of $390\,\mu$s and $2838\,\mu$s, without a matching
global RMSE gain. On the RMSE--CPU plane, EKF, RA-EKF, SOGI-FLL, SOGI-PLL, and
ZCD are Pareto-optimal; in practice, the protection-relevant subset is
\{SOGI-FLL, SOGI-PLL, EKF, RA-EKF\}, since ZCD reaches the front only by
near-zero cost and not by usable composite-stress accuracy.

For deployment screening within this benchmark, estimator choice remains tied
to disturbance family, operating objective, and compute budget. SOGI-PLL is
strongest on ramps, EKF on magnitude steps and out-of-band interference,
RA-EKF on modulation, and PLL on the composite IBR event. Figure~\ref{fig:dashboard_mega}(g)
keeps SOGI-FLL, SOGI-PLL, EKF, and RA-EKF in the main low-cost cluster below
$2\,\mu$s per sample, whereas TFT and IpDFT move rightward in cost but improve
trip-risk behavior in the composite cases. Those are conditional results, not a
universal hierarchy.

\begin{figure*}[!t]
  \centering
  \includegraphics[width=\textwidth,height=0.80\textheight,keepaspectratio]{Fig2_Mega_Dashboard.pdf}
  \caption{Representative benchmark views. Panels (a)--(d) show single-scenario
    responses summarized over the same 60 Monte Carlo runs used throughout the
    study. Panel (e) gives a disturbance scorecard across benchmark families.
    Panel (f) shows the scenario-level RMSE distribution grouped by estimator
    family. Panels (g) and (h) show global cost-versus-accuracy and
    cost-versus-trip-risk views; the dashed vertical line in both panels marks
    the 100-$\mu$s/sample real-time budget implied by the 10-kHz acquisition
    rate. The dashboard shows that estimator selection depends on disturbance
    family, risk metric, and compute budget.}
  \label{fig:dashboard_mega}
\end{figure*}

\section{Conclusion}
\label{sec:conclusion}

This paper compares 16 frequency estimators over 34 scenarios in a local
relay-class measurement setting that combines standard dynamic tests with
IBR-specific composite stress. The results show disturbance dependence:
estimator rankings move with both the disturbance family and the metric used to
read them. Standard compliance-style tests remain necessary reference cases,
but they are not sufficient by themselves for protection-oriented evaluation in
IBR-dominated grids.

We highlight three findings. First, the latency--robustness trade-off is
structural: nine estimators win at least one RMSE scenario, and no estimator
wins more than 10 of the 34. Windowed spectral methods reduce steady harmonic
error but pay observation delay, while low-cost loop-based methods remain
attractive under tight computational budgets. Second, Kalman-family methods,
especially EKF and RA-EKF in this dataset, remain competitive across the
widest range of families. RA-EKF is evaluated here as one Kalman-family
variant rather than as a new estimator family. Third, the hardest cases are
the IBR multi-event sequence, the $60^\circ$ phase-jump cases, and the noisy
ringdown variants; these are also the cases in which RMSE, peak error, and
trip-risk select different estimators.

The study is limited to the scenario set and relative computational-cost
measurements reported here. CPU time is used as a software-level cost indicator,
not as hardware latency, and the study remains simulation-only. Hardware-in-the-loop
validation, wider coverage of converter-control interactions, and sensitivity
analysis of the tuning protocol are the natural next steps.

From a deployment standpoint, the benchmark identifies a small practical core.
SOGI-PLL is the strongest low-cost choice for ramp-dominated operation. EKF is
the most consistent model-based baseline for magnitude steps, out-of-band
interference, and balanced composite-event behavior. RA-EKF earns its place not
as a universal winner but as a higher-cost Pareto point with the best mean RMSE
and the strongest modulation-family results. IpDFT and TFT remain relevant when
limiting trip-risk matters more than minimizing average error, provided their
observation delay is acceptable.

In practice, the benchmark does not support a single default estimator for all
low-inertia IBR conditions. It supports a screening workflow: narrow
candidates by disturbance family and compute budget, then resolve the final
choice with the metric aligned to the protection objective. In that workflow,
ramps mostly test tracking bandwidth, jumps test recovery after
discontinuities, and harmonic-rich cases test robustness to leakage and
interharmonics.

\ifCLASSOPTIONcaptionsoff
  \newpage
\fi

% Generated by IEEEtran.bst, version: 1.14 (2015/08/26)
\begin{thebibliography}{10}
\providecommand{\url}[1]{#1}
\csname url@samestyle\endcsname
\providecommand{\newblock}{\relax}
\providecommand{\bibinfo}[2]{#2}
\providecommand{\BIBentrySTDinterwordspacing}{\spaceskip=0pt\relax}
\providecommand{\BIBentryALTinterwordstretchfactor}{4}
\providecommand{\BIBentryALTinterwordspacing}{\spaceskip=\fontdimen2\font plus
\BIBentryALTinterwordstretchfactor\fontdimen3\font minus
  \fontdimen4\font\relax}
\providecommand{\BIBforeignlanguage}[2]{{%
\expandafter\ifx\csname l@#1\endcsname\relax
\typeout{** WARNING: IEEEtran.bst: No hyphenation pattern has been}%
\typeout{** loaded for the language `#1'. Using the pattern for}%
\typeout{** the default language instead.}%
\else
\language=\csname l@#1\endcsname
\fi
#2}}
\providecommand{\BIBdecl}{\relax}
\BIBdecl

\bibitem{Milano2018LowInertia}
F.~Milano, F.~D{\"o}rfler, G.~Hug, D.~J. Hill, and G.~Verbi{\v{c}},
  ``Foundations and challenges of low-inertia systems (invited paper),'' in
  \emph{2018 Power Systems Computation Conference (PSCC)}, 2018, pp. 1--25.

\bibitem{Pinheiro2025}
B.~Pinheiro, M.~C. Llerena~Velasquez, G.~Faria, A.~F.~C. Aquino, D.~Issicaba,
  and D.~Dotta, ``Evaluation of the impact of {IBR} on the frequency dynamics
  in the {Brazilian Power System},'' \emph{Sustainable Energy, Grids and
  Networks}, vol.~43, p. 101827, 2025.

\bibitem{Ponce2025}
I.~Ponce, C.~Rahmann, and P.~Pourbeik, ``Identification of weak frequency areas
  in low inertia systems,'' \emph{IEEE Access}, vol.~13, pp. 90\,262--90\,272,
  2025.

\bibitem{Bernal2025}
R.~Bernal and F.~Milano, ``A complex frequency-based control for inverter-based
  resources,'' \emph{Journal of Modern Power Systems and Clean Energy},
  vol.~13, no.~5, pp. 1630--1641, 2025.

\bibitem{Zoghby2024}
P.~Zoghby, B.~Marinescu, and A.~Ross{\'e}, ``A new benchmarking method for the
  study of dynamic interactions in island power systems with hybrid
  conventional and inverter-based resources controlled in grid-forming mode,''
  in \emph{2024 International Conference on Renewable Energies and Smart
  Technologies (REST)}, Prishtina, Kosovo, 2024.

\bibitem{Belega2014}
D.~Belega and D.~Dallet, ``An improved {DFT}-based method for the accurate
  measurement of frequency, amplitude, and phase of power system signals,''
  \emph{IEEE Transactions on Instrumentation and Measurement}, vol.~63, no.~2,
  pp. 364--371, 2014.

\bibitem{PlatasGarza2011}
F.~J. Platas-Garc{\'i}a and J.~A. de~la O~Serna, ``Dynamic phasor estimation
  using a {Taylor--Fourier} multisignal model,'' \emph{IEEE Transactions on
  Instrumentation and Measurement}, vol.~60, no.~6, pp. 2112--2121, 2011.

\bibitem{Golestan2017}
S.~Golestan, J.~M. Guerrero, and J.~C. Vasquez, ``Three-phase {PLL}s: A review
  of recent advances,'' \emph{IEEE Transactions on Power Electronics}, vol.~32,
  no.~3, pp. 1894--1907, 2017.

\bibitem{Wright2019}
P.~S. Wright, P.~N. Davis, K.~Johnstone, G.~Rietveld, and A.~J. Roscoe, ``Field
  measurement of frequency and {RoCoF} in the presence of phase steps,''
  \emph{IEEE Transactions on Instrumentation and Measurement}, vol.~68, no.~6,
  pp. 1688--1695, 2019.

\bibitem{Ferrero2020EKF}
R.~Ferrero, P.~A. Pegoraro, and S.~Toscani, ``Dynamic synchrophasor estimation
  by extended {Kalman} filter,'' \emph{IEEE Transactions on Instrumentation and
  Measurement}, vol.~69, no.~7, pp. 4818--4826, 2020.

\bibitem{Chamorro2021}
H.~R. Chamorro, A.-J. Guel-Cortez, Y.~Africano, F.~Jimenez, R.~Tonkoski, and
  W.~Martinez, ``Grid frequency estimation based on extended {Kalman} filter in
  connected {AC} microgrids,'' in \emph{2021 IEEE 30th Int. Symp. on Industrial
  Electronics (ISIE)}, 2021, pp. 1--6.

\bibitem{Netto2018KoopmanKF}
M.~Netto and L.~Mili, ``A robust data-driven {Koopman Kalman} filter for power
  systems dynamic state estimation,'' \emph{IEEE Transactions on Power
  Systems}, vol.~33, no.~6, pp. 7228--7237, 2018.

\bibitem{Ting2022}
J.~Ting, B.~Wang, R.~W. Kenyon, and A.~Hoke, ``Evaluating methods for measuring
  grid frequency in low-inertia power systems,'' in \emph{2022 IEEE Kansas
  Power and Energy Conference (KPEC)}, 2022, {NREL}/CP-6A40-83784.

\bibitem{IEEE_C37_118}
{IEEE Power and Energy Society}, ``{IEEE} standard for synchrophasor
  measurements for power systems ({C37.118.1-2011}, amendment
  {C37.118.1a-2014}),'' IEEE Std C37.118.1, 2011.

\bibitem{IEC_IEEE_60255_118_1_2018}
{IEC and IEEE}, ``{IEC/IEEE} international standard 60255-118-1: Synchrophasor
  measurements for power systems,'' IEC/IEEE 60255-118-1:2018, 2018.

\end{thebibliography}

\end{document}
